\documentclass[instructions.tex]{subfiles}
\begin{document}

\chapter{Du châtiment du péché.}


Par la même raison que Dieu hait le péché, il doit le punir; car le péché, dit saint Augustin, ne serait plus péché, s'il ne devait pas être puni\footnote{Peccatum non esset, si puniendum non esset. S. Aug.}. Dieu agit toujours en Dieu; il doit aussi, en punissant le péché, le punir en Dieu. Il le punit en cette vie et en l'autre, il le punit sans acception de personne, et toujours avec une souveraine justice.

\primo En cette vie. A peine Adam eut-il péché, qu'il fut puni, dépouillé de ses privilèges, et tous ses descendants enveloppés dans son malheur. Les fléaux, les misères humaines et la mort sont les suites de ce péché; de telle sorte que, tandis qu'il y aura sur la terre une goutte de sang de ce premier homme, Dieu sera irrité contre ses descendants. Tel est le châtiment du péché de notre premier père, châtiment plein d'équité et digne de Dieu. Si un roi qui punit un crime de lèse-majesté a droit d'envelopper les enfants dans la disgrâce d'un père coupable, Dieu, à plus forte raison, le peut avec justice.

Combien d'autres châtiments du péché voyons-nous dans l'histoire sainte! Caïn maudit de Dieu, proscrit de l'univers; le monde noyé dans un déluge, pour les impuretés de ses habitants; une pluie de soufre et de feu sur quatre villes impudiques; des serpens de feu envoyés aux Juifs murmurateurs; le roi Saül, la reine Jésabel, le détestable Antiochus, et tant d'autres, punis sévèrement de Dieu; l'armée d'Israël taillée en pièces pour le vol du seul Acham; cinquante mille habitants de Betsamée frappés de mort pour avoir profané l'arche sainte; la nation des Juifs aujourd'hui errante, en exécration à tous les peuples, visiblement punie de Dieu pour avoir méprisé ses grâces et sa parole. Combien, hélas! d'autres exemples!

Les fléaux du pécheur, dit l'Écriture, sont en grand nombre. Si nous éprouvons les malheurs des temps, les guerres, les mortalités, le dérangement des saisons, c'est parce que les crimes des hommes, les larcins, les impudicités, les blasphèmes, les sacrilèges et les impiétés inondent la terre. Si vous violez ma loi, disait Dieu à son peuple, je vous réduirai aussitôt à la pauvreté\footnote{Visitabo vos velociter in egestate. Lev.26.}. Si vous continuez dans vos prévarications, je redoublerai vos maux sept fois autant\footnote{Addam plagas vestras in septuplum. Ibid.}. Le ciel sera pour vous de fer, et la terre, de bronze et d'airain\footnote{Dabo vobis cœlum desuper sicut ferrum, et terram æneam. Ibid.}.

C'est donc dans nous-mêmes que nous devons chercher la cause ordinaire de nos maux; c'est aussi dans nous que nous en trouverons le remède par une sincère conversion. Notre retour à Dieu désarmera sa colère. En exerçant sa justice ici-bas, il n'oublie point sa miséricorde. Il frappe les pécheurs pour les rappeler à lui, et les justes pour les purifier. Que ceux-là donc sont malheureux, qui s'endurcissent sous les coups de sa main puissante!

Les autres, qui n'éprouvent aucune disgrâce, en vivant dans le désordre, ne sont pas moins malheureux. Une grande prospérité, jointe à une vie criminelle, est un grand châtiment de Dieu. S'il épargne certains scélérats en ce monde, c'est pour les punir en l'autre avec plus de sévérité.

\secundo En l'autre vie. Jugeons de la haine de Dieu pour le péché, par les châtiments de l'autre vie. Quand je vois un réprouvé poussant des cris lamentables, condamné à brûler dans des feux dévorans pendant des siècles sans fin; quand je vois des âmes saintes expier en purgatoire des fautes légères par des tourmens inexprimables, je suis saisi d'étonnement et de frayeur, je m'écrie: O que Dieu est terrible dans ses conseils\footnote{Terribilis in consiliis.}!

Mais d'un autre côté, quand je pense à la grandeur de Dieu et à la bassesse de l'homme, quand je considère qu'un Dieu a été traité avec mépris par sa créature, que cette créature a abusé de la bonté de son Créateur, qui voulait la sauver; qu'elle l'a affligé, qu'elle lui a refusé l'hommage et la soumission, je ne m'étonne plus d'un châtiment si rigoureux.

Je comprends qu'ayant négligé de mériter un bonheur qui lui était destiné, elle doit en être privée pour toujours. Je comprends qu'un néant, qui a osé se révolter contre l'être souverain, doit être puni souverainement, et qu'ayant eu la malice de mettre son Dieu en comparaison avec des plaisirs et des objets frivoles, ayant même préféré ces plaisirs et ces objets à la possession d'un bien infini et éternel, je comprends qu'il est juste que, pour le punir, Dieu se serve de la même mesure dont l'homme s'est servi pour l'offenser, et que cette offense soit mesurée et punie par une peine éternelle et sans fin; et je m'écrie: Vous êtes juste, Seigneur, vos jugemens sont pleins d'équité\footnote{Justus es, Domine, et rectum judicium tuum.}.

\tertio Il punit de la sorte sans acception de personne. Il faut que le péché soit bien horrible, et infiniment plus injurieux à Dieu que nous ne pouvons le comprendre, puisque Dieu, qui n'est que bonté et miséricorde, le punit avec tant de sévérité. Si un père, qui aime tendrement son fils, le condamnait à périr par le feu; et si, après l'avoir fait jeter dans une fournaise ardente, ce père insultait à l'infortune de cet enfant malheureux, en lui disant: Je fais ma gloire de ton supplice; les cris et les larmes ne me touchent point; tu n'es plus mon fils, ne me regarde plus comme ton père; vous jugeriez que ce misérable enfant aurait commis un énorme attentat.

Jugez par proportion de l'attentat commis contre Dieu, et de la haine que Dieu a pour le péché. Dieu aime l'homme infiniment plus que le père le plus tendre n'aime son fils; cependant, quand il s'agit de punir le péché mortel, il précipiterait non-seulement un homme, mais une infinité d'hommes sans aucune acception, dans des abîmes de feu, et mettrait sa gloire à les punir pendant toute l'éternité.

Dieu aimait les anges, ils étaient ses plus nobles créatures; néanmoins, pour un seul péché, il en précipite des millions sans miséricorde et les réprouve. Abraham, Moïse, Élie, étaient les amis de Dieu; cependant ces grands saints, ni aucun homme juste, n'eussent jamais été sauvés, s'ils fussent morts souillés d'un seul péché mortel. O mon Dieu, que le péché est donc un grand mal, puisque vous le punissez avec tant de rigueur! Quel sujet de trembler, puisque l'homme le plus saint, s'il vient à tomber et à mourir dans un seul péché mortel, est perdu sans ressource! Que nous sommes aveugles si nous ne le craignons pas, si nous en aimons l'occasion, et si nous vivons sans vigilance et sans précautions pour nous-mêmes!


\section{Résolutions}

Je vais commencer à me punir moi-même des péchés de ma vie passée.

\primo Pour cela, je ferai tous les jours quelques actes de contrition.

\secundo Et je vengerai Dieu sur tout ce qui a péché en moi: sur mon esprit, en réprimant les pensées vaines et inutiles; sur mon cœur, en le remplissant de sentimens de douleur pour mes transgressions; et sur toutes mes facultés corporelles, en leur imposant des privations utiles.

\prayer{Mon Dieu, armez-moi d'une sainte haine contre moi-même.}

\end{document}
